\begin{table*}[t]
\centering
\caption{ARA 的三个版本及证据归属。提交简称只定位源码，不保证重建精确运行工作树。}
\label{tab:ara-version-comparison}
\footnotesize
\setlength{\tabcolsep}{4pt}
\renewcommand{\arraystretch}{1.15}
\begin{tabularx}{\textwidth}{@{}p{24mm}Y Y Y@{}}
\toprule
对象 & 上游 ARA 原型 & 本地 point-v1（CARA） & 本地 trajectory-v2 \\
\midrule
版本依据 & 固定原型 \code{edc3b123456c} & 运行记录 \code{cd2977a}；部署修复后对应 \code{868ca73} & 审计源码 \code{2871bc1} 与 v2 模板 \\
局部样本 & 冻结末位置 I/O & 冻结末位置 I/O；每类 64 条 & 基座多位置轨迹；模板每类 96 条、至多八个位置 \\
目标与参数化 & 硬 $k$NN、负推远距离；全矩阵或低秩 & 归一化软邻域、非负间隔惩罚、归一化 Gram 项；秩至多 128 & 同步位软邻域、提示内衰减权重；未按初始值归一化的 Gram 项 \\
恢复与部署 & 低秩仅清零 $B$；可有行归一化目标/部署差异 & 恢复全部 $A,B$；无行归一化；QR/SVD 输出检查 & 轨迹目标、部署增益及漂移保护；分阶段验收 \\
搜索与评分 & 默认 200 次；历史分段评分 & 六维、36+84 次；原始 Keywords 与 KL & 八维、8+24+88 次；连续前缀分数与 KL \\
允许的证据 & E1 源码与 E2 数学分析 & E6 一次 120 次试验搜索；效果门槛失败 & E1 已实现控制；无指定归档中的效果结果 \\
\bottomrule
\end{tabularx}
\end{table*}
